%!TEX root = ED_Lista_Fila_Pilha

\section{Pilha}
\begin{frame}[label=notas]
	\begin{itemize}
		\item Lista LIFO (Last In, First Out)
		\item Os elementos são colocados na pilha e retirados (ou processados) em ordem inversa a de chegada
			\begin{itemize}
				\item Inserção \alert<1>{apenas no topo} da pilha
				\item Remoção \alert<1>{apenas do topo} da pilha
			\end{itemize}
	\end{itemize}
	\begin{center}
		\includegraphics[width=.5\textwidth]{stack}
		%source: http://commons.wikimedia.org/wiki/File:Data_Queue.svg
	\end{center}
\end{frame}

\begin{frame}[label=notas]
	Implementação? 
	\\\pause\vfill
	Operações possíveis:
	\begin{itemize}
		\item limpar a pilha 
		\item verificar se a pilha está vazia 
		\item colocar elemento no topo da pilha 
		\item retirar o elemento no topo da pilha
		\item etc.
	\end{itemize}
\end{frame}

% \lstset{morekeywords={stack_node_t, stack_t}}
% \begin{frame}[fragile,label=notas]
% 	\vspace{-2\baselineskip}
% 	\begin{center}
% 		\begin{minipage}{1.02\textwidth}
% 			\begin{lstlisting}[title=Criar elemento]
% 				stack_node_t *stack_new_node(data_t item, stack_node_t *next){
% 				  stack_node_t *stack_item = (stack_node_t *) malloc(sizeof(stack_node_t));
% 				  if (stack_item) {
% 				    stack_item->item = item;
% 				    stack_item->next = next;
% 				  }
% 				  return stack_item;
% 				}
% 			\end{lstlisting}
% 		\end{minipage}
% 	\end{center}
% 	\vspace{-2em}\pause
% 	\begin{center}
% 		\begin{minipage}{.8\textwidth}
% 			\begin{lstlisting}[title=Criar pilha]
% 			                	stack_t *stack_new(data_t *item) {
% 			                	  stack_t *pilha = (stack_t *)malloc(sizeof(stack_t));
% 			                	  if(item) push(pilha, *item);
% 			                	  else pilha->first = NULL;
% 			                	  return pilha;
% 			                	}
% 			\end{lstlisting}			
% 		\end{minipage}
% 	\end{center}
% \end{frame}
% \lstset{morekeywords={stack_new_node, stack_new}}

% \begin{frame}[fragile,label=notas]
% 	\begin{center}
% 		\begin{minipage}{.75\textwidth}
% 			\begin{lstlisting}[title=Inserir]
% 				void push(stack_t *pilha, data_t item) {
% 				  pilha->first = stack_new_node(item, pilha->first);
% 				}
% 			\end{lstlisting}			
% 		\end{minipage}
% 	\end{center}
% 	\pause
% 	\begin{center}
% 		\begin{minipage}{.5\textwidth}
% 			\begin{lstlisting}[title=Remover]
% 				stack_node_t *pop(stack_t *pilha) {
% 				  stack_node_t *aux = pilha->first;
% 				  pilha->first = aux->next;
% 				  return aux;
% 				}
% 			\end{lstlisting}			
% 		\end{minipage}
% 	\end{center}
% \end{frame}
% \lstset{morekeywords={push, pop}}

% \begin{frame}[fragile,label=notas]
% 	\begin{center}
% 		\begin{minipage}{.5\textwidth}
% 			\begin{lstlisting}[title=Esvaziar pilha]
% 				void stack_empty(stack_t *pilha) {
% 				  while (pilha->first)
% 				    free(pop(pilha));
% 				  pilha->first = NULL;
% 				}
% 			\end{lstlisting}			
% 		\end{minipage}
% 	\end{center}
% 	\pause
% 	\begin{center}
% 		\begin{minipage}{.6\textwidth}
% 			\begin{lstlisting}[title=Verificar se está vazia]
% 				int stack_is_empty(stack_t *pilha) {
% 				  return (pilha->first == NULL);
% 				}
% 			\end{lstlisting}			
% 		\end{minipage}
% 	\end{center}
% \end{frame}
% \lstset{morekeywords={stack_empty, stack_is_empty}}


\begin{frame}
	É ideal para processamento de estruturas aninhadas de profundidade imprevisível.
	\\\pause\vfill
	Uma pilha contém uma sequência de  obrigações adiadas. A ordem de remoção garante que as  estruturas mais internas serão processadas antes das mais externas.
		\begin{itemize}
			\item Quando é necessário caminhar em um conjunto de dados e guardar uma lista de coisas a fazer posteriormente.
			\item O controle de sequências de chamadas de subprogramas.
			\item A sintaxe de expressões aritméticas.
		\end{itemize}
	\pause\vfill
	As pilhas ocorrem em estruturas de natureza recursiva (como árvores) e  são utilizadas para implementar a recursividade.
\end{frame}

\begin{frame}
	\framesubtitle{\subsecname\ - Aplicação}
	Controle de delimitadores em uma equação (rastrear o escopo).
	\subitem[]{$\{(A+B)*C+(D+E)/[F+G]+H\}/I$}
	\pause\vfill
	\begin{enumerate}
		\item se um inicializador de escopo for encontrado, ele é empilhado
		\item se um finalizador de descopo for encontrado, a pilha é verificada
			\begin{itemize}
				\item se estiver vazia $\rightarrow$ equação incorreta
				\item se não, desempilhar e comparar ao finalizador
			\end{itemize}
	\end{enumerate}
	\vfill
	Ao final, a pilha deve estar vazia.
\end{frame}


\begin{frame}[fragile]
	Aplicações de Pilhas: \alert<1>{recursividade}\pause
		\begin{itemize}
			\item a solução de um problema depende da solução de instâncias menores do mesmo problema
			\item um programa de computador finito para processar algo potencialmente infinito
		\end{itemize}
		\pause
		\begin{columns}[T]
		\column{.5\textwidth}
			\begin{lstlisting}[title=Fatorial iterativo]
				int fatorial(int n) {
				  int result = 1;
				  while(n > 1) {
				    result *= n;
				    n -= 1;
				   }
				   return result;
				 }
			\end{lstlisting}
			\pause
		\column{.45\textwidth}
			\begin{lstlisting}[title=Fatorial recursivo]
			                	int fatorial(int n) {
			                	  if(n<=1)
			                	    return 1;
			                	  else
			                	    return n*fatorial(n-1);
			                	 }
			\end{lstlisting}			
	\end{columns}
\end{frame}

\begin{frame}
	Quanto é $4!$ ?\pause
	\begin{columns}
		\column{.6\textwidth}
			\begin{figure}
				\centering
				\begin{tikzpicture}[node distance=1em,minimum width=15em, minimum height=1.8em]
					\tikzstyle{blockEdgePortion} = [black];
					\tikzstyle{blockEdge} = [blockEdgePortion,->];
								
					\node[block] (f4) {factorial(4) $\leftarrow$ 4*fatorial(3)};
					\node<3->[block, above=of f4] (f3) {fatorial(3) $\leftarrow$ 3*fatorial(2)} edge[line,<-] (f4);
					\node<4->[block, above=of f3] (f2) {fatorial(2) $\leftarrow$ 2*fatorial(1)} edge[line,<-] (f3);
					\node<5->[block, above=of f2] (f1) {fatorial(1) $\leftarrow$ 1} edge[line,<-] (f2);
				\end{tikzpicture}
			\end{figure}
		\column{.35\textwidth}
			\uncover<6>{
				\begin{figure}
					\centering
					\begin{tikzpicture}[node distance=1em,minimum width=5em, minimum height=1.8em, color=blueUnB]
						\node[block] (f1r) {1};
						\node[block, below=of f1r] (f2r) {2*1 = 2} edge[line,<-] (f1r);
						\node[block, below=of f2r] (f3r) {3*2 = 6} edge[line,<-] (f2r);
						\node[block, below=of f3r] (f4r) {4*6 = \alert{24}} edge[line,<-] (f3r);
					\end{tikzpicture}
				\end{figure}
			}
	\end{columns}
\end{frame}